\chapter{The Strategic Opportunity: Problem \& Objectives}
%\label{sec:excellence}

\section{The Core Problem}
Travelers, from leisure visitors to business professionals, face significant and costly friction in planning and executing urban itineraries, which stems from a core problem of fragmented data. Critical decision-making data on public transit schedules (GTFS), ride-hail fares, live traffic, and operating hours exists in disconnected silos.
\\
This forces users to manually compare multiple apps, leading to:
\begin{itemize}
    \item  Inefficient planning under complex constraints (time, budget, transfers).
    \item Suboptimal in-the-moment transport choices, failing to balance cost, ETA.
    \item Reactive, unactionable responses to congestion and crowding.
\end{itemize}

\section{PoC Objectives}
Success is defined by achieving statistically significant signals against predefined goals: 
\begin{itemize}
    \item \textbf{Feasibility:} $\geq90\%$ of all generated test itineraries must successfully satisfy all hard constraints (time windows, budgets, mobility limits).
    \item \textbf{Adoption:} $\geq60\%$ user acceptance of the system’s top-3 generated daily plans.
    \item \textbf{In-Trip Value:} Achieve an on-time arrival rate of $\geq80\%$ and a missed-connection rate of $\leq5\%$.
    \item \textbf{Safety Value:} Deliver alert precision of $\geq80\%$ and limit user exposure to congested segments to $\leq15$ minutes per day.
    \item \textbf{Qualitative:} Secure an average CSAT score of $\geq4.2/5$ on both planning and in-trip guidance
\end{itemize}



%\cleardoublepage

\glsresetall



